\subsection{Associative Algebra Structures}

% Role: algebraic envelope of endomorphism algebras, polynomial actions, and
% finite-dimensional division algebra examples.

The module viewpoint in Definition \ref{def:operator-module-and-matrix-module} uses one operator at a time. An algebra keeps track of a whole vector space of operators or products at once. The guiding example is $\LL(\mathcal{V})$, where addition and scalar multiplication are pointwise and multiplication is composition.

\begin{definition} \label{def:unital-associative-algebra} \hlt{Unital Associative Algebra}\\
Let $\mathbb{F}$ be a field.
\begin{enumerate}[label=(\roman*)]
\setlength{\itemsep}{0pt}
\item An \hlt{associative $\mathbb{F}$-algebra} is an $\mathbb{F}$-vector space $\mathcal{A}$ and unital ring $(\mathcal{A},+,\cdot,1_{\mathcal{A}})$ such that $r(ab)=(ra)b=a(rb)$ for all $r\in\mathbb{F}$ and $a,b\in\mathcal{A}$.
\item $\mathcal{A}$ is \hlt{finite-dimensional} iff $\dim_{\mathbb{F}}\mathcal{A}<\infty$, and \hlt{commutative} iff $ab=ba$ for all $a,b\in\mathcal{A}$.
\item $a\in\mathcal{A}$ is \hlt{invertible} iff $\exists b\in\mathcal{A}$ such that $ab=ba=1_{\mathcal{A}}$.
\item In this subsection, \hlt{algebra} means unital associative $\mathbb{F}$-algebra unless stated otherwise.
\end{enumerate}
\end{definition}

\begin{remark} \label{rem:algebra-convention-relative-to-ring-definition} \hlt{Convention}\\
Relative to Definition \ref{def:algebra-over-field}, assume multiplication is associative and unital.\\
Thus, \hlt{algebra} means unital associative $\mathbb{F}$-algebra.
\end{remark}

\begin{definition} \label{def:center-and-central-algebra} \hlt{Center and Central Algebras}
\begin{enumerate}[label=(\roman*)]
\setlength{\itemsep}{0pt}
\item The \hlt{center} of an $\mathbb{F}$-algebra $\mathcal{A}$ is $Z(\mathcal{A})\defeq\{a\in\mathcal{A}: ax=xa\ \forall x\in\mathcal{A}\}$.
\item $\mathcal{A}$ is \hlt{central over $\mathbb{F}$} iff $Z(\mathcal{A})=\mathbb{F}1_{\mathcal{A}}\defeq\{r1_{\mathcal{A}}:r\in\mathbb{F}\}$.
\end{enumerate}
\end{definition}

\begin{proposition} \label{prop:field-embeds-in-center-of-algebra} \hlt{The Scalar Field Sits in the Center}\\
If $\mathcal{A}$ is an $\mathbb{F}$-algebra with $1_{\mathcal{A}}\neq0$, then $\iota_{\mathcal{A}}:\mathbb{F}\hookrightarrow Z(\mathcal{A}),\quad r\mapsto r1_{\mathcal{A}}$, is an injective unital ring homomorphism.\\
Conversely, if $R$ is a unital ring and $\mathbb{F}\leq Z(R)$ is a subfield, then $R$ is an $\mathbb{F}$-algebra via $r\cdot x\defeq rx$.
\begin{proof}
\begin{enumerate}[label=\roman*.]
\setlength{\itemsep}{0pt}
\item[$(\Rightarrow)$] Write $\iota=\iota_{\mathcal{A}}$. For $r,s\in\mathbb{F}$, $\iota(r+s)=(r+s)1_{\mathcal{A}}=\iota(r)+\iota(s)$, $\iota(rs)=(rs)1_{\mathcal{A}}=(r1_{\mathcal{A}})(s1_{\mathcal{A}})$, $\iota(1)=1_{\mathcal{A}}$.\\
For $r\in\mathbb{F}$ and $a\in\mathcal{A}$, $\iota(r)a=(r1_{\mathcal{A}})a=r(1_{\mathcal{A}}a)=ra=r(a1_{\mathcal{A}})=a(r1_{\mathcal{A}})=a\iota(r)$, so $\operatorname{im}\iota\subseteq Z(\mathcal{A})$.\\
If $\iota(r)=0$ and $r\neq0$, then $1_{\mathcal{A}}=\iota(r^{-1})\iota(r)=0$, contradiction; hence $\ker\iota=\{0\}$.
\item[$(\Leftarrow)$] Define $r\cdot x=rx$. The field laws in $\mathbb{F}\subseteq R$ give the vector-space axioms, and $1_{\mathbb{F}}\cdot x=x$.\\
For $r\in\mathbb{F}$ and $x,y\in R$, centrality gives $r\cdot(xy)=rxy=(rx)y=x(ry)=x(r\cdot y)$, so the scalar action is compatible with the ring product.
\end{enumerate}
\hlt{Structural Commentary}: An algebra structure is a choice of central scalar copy $\mathbb{F}1_{\mathcal{A}}\subseteq Z(\mathcal{A})$.\\
\hlt{Key Insight}: Scalars are precisely central multipliers of the unit.
\end{proof}
\end{proposition}

\begin{method} \label{meth:structure-constants-algebra} \hlt{Building Multiplication from Structure Constants}\\
Let $V=\bigoplus_{i\in I}\mathbb{F}b_i$. Choose $c_{ij}^k\in\mathbb{F}$ with $\{k:c_{ij}^k\neq0\}$ finite for each $(i,j)$.
\begin{enumerate}[label=(\roman*)]
\setlength{\itemsep}{0pt}
\item Define $b_i b_j\defeq\sum_{k\in I}c_{ij}^k b_k$ and extend bilinearly: $(\sum_i r_i b_i)(\sum_j s_j b_j)\defeq\sum_{i,j,k}r_i s_j c_{ij}^k b_k$.
\item The product is associative iff $\sum_{\ell\in I}c_{ij}^{\ell}c_{\ell k}^{m}=\sum_{\ell\in I}c_{jk}^{\ell}c_{i\ell}^{m}$ for all $i,j,k,m\in I$.
\item The product is commutative iff $c_{ij}^k=c_{ji}^k$ for all $i,j,k\in I$.
\item A basis vector $b_e$ is an identity iff $c_{ej}^m=\delta_j^m=c_{je}^m$ for all $j,m\in I$.
\end{enumerate}
\end{method}

\begin{example} \label{ex:basic-associative-algebras} \hlt{Basic Algebras}
\begin{enumerate}[label=(\roman*)]
\setlength{\itemsep}{0pt}
\item If $\mathbb{F}\subseteq\mathbb{E}$ are fields, then $\mathbb{E}$ is an $\mathbb{F}$-algebra by $r\cdot x=rx$ and multiplication in $\mathbb{E}$.
\item $\mathbb{F}[x]$ is a commutative $\mathbb{F}$-algebra with unit $1$ and scalar embedding $r\mapsto r$.
\item $\operatorname{Mat}(n,n,\mathbb{F})$ is an $\mathbb{F}$-algebra with unit $I_n$ and product $AB$.
\item For an $\mathbb{F}$-vector space $\mathcal{V}$, $\LL(\mathcal{V})=\operatorname{End}_{\mathbb{F}}(\mathcal{V})$ is an $\mathbb{F}$-algebra with unit $I_{\mathcal{V}}$ and product $S\circ T$.
\end{enumerate}
\end{example}

\begin{proposition} \label{prop:endomorphism-algebra-matrix-algebra-isomorphism} \hlt{Endomorphisms and Matrices Are the Same Algebra after Choosing a Basis}\\
Let $\dim\mathcal{V}=n$ and fix an ordered basis $\mathcal{B}$ of
$\mathcal{V}$.\\
Then $\mu_{\mathcal{B}}:\LL(\mathcal{V})\to\operatorname{Mat}(n,n,\mathbb{F})$, $T\mapsto[T]_{\mathcal{B}}$, is an $\mathbb{F}$-algebra isomorphism.
\begin{proof}
By Theorem \ref{thm:linear-maps-and-matrices-same-data}, $\mu_{\mathcal{B}}$ is an $\mathbb{F}$-linear bijection.\\
For $S,T\in\LL(\mathcal{V})$ and $x\in\mathcal{V}$, $[STx]_{\mathcal{B}}=[S]_{\mathcal{B}}[Tx]_{\mathcal{B}}=[S]_{\mathcal{B}}[T]_{\mathcal{B}}[x]_{\mathcal{B}}$, hence $\mu_{\mathcal{B}}(ST)=\mu_{\mathcal{B}}(S)\mu_{\mathcal{B}}(T)$.\\
Also $\mu_{\mathcal{B}}(I_{\mathcal{V}})=I_n$; therefore $\mu_{\mathcal{B}}$ is an $\mathbb{F}$-algebra isomorphism.\\
\hlt{Structural Commentary}: A basis is a coordinate bridge from abstract endomorphisms to concrete matrices.\\
\hlt{Key Insight}: Composition becomes matrix multiplication after taking coordinate columns.
\end{proof}
\end{proposition}
